\hypertarget{appendix-u-the-standard-library-concurrency-toolkit-a-cppreference-breakdown}{%
\section{Appendix U: THE STANDARD LIBRARY CONCURRENCY TOOLKIT (A
Cppreference
Breakdown)}\label{appendix-u-the-standard-library-concurrency-toolkit-a-cppreference-breakdown}}

If you look at the \passthrough{\lstinline!<thread>!} or
\passthrough{\lstinline!<atomic>!} pages on cppreference, they are
written in ``Standardese'' (the language of the ISO C++ committee). This
appendix translates the most critical concurrency tools into ``Head
First'' English.

\hypertarget{u.1-thread-and-jthread}{%
\subsection{\texorpdfstring{U.1 \texttt{\textless{}thread\textgreater{}}
and
\texttt{\textless{}jthread\textgreater{}}}{U.1 \textless thread\textgreater{} and \textless jthread\textgreater{}}}\label{u.1-thread-and-jthread}}

\hypertarget{stdthreadhardware_concurrency}{%
\subsubsection{\texorpdfstring{\texttt{std::thread::hardware\_concurrency()}}{std::thread::hardware\_concurrency()}}\label{stdthreadhardware_concurrency}}

\begin{itemize}
\tightlist
\item
  \textbf{Cppreference says}: Returns the number of concurrent threads
  supported by the implementation.
\item
  \textbf{Head First Translation}: ``How many physical/logical CPU cores
  do I have?''
\item
  \textbf{Godhood Tip}: Do not spawn 1,000 threads if you only have 8
  cores. The OS will spend all its time context-switching between
  threads instead of actually doing work. Create a Thread Pool with
  exactly \passthrough{\lstinline!hardware\_concurrency()!} workers.
\end{itemize}

\hypertarget{stdthis_threadyield}{%
\subsubsection{\texorpdfstring{\texttt{std::this\_thread::yield()}}{std::this\_thread::yield()}}\label{stdthis_threadyield}}

\begin{itemize}
\tightlist
\item
  \textbf{Cppreference says}: Provides a hint to the implementation to
  reschedule the execution of threads, allowing other threads to run.
\item
  \textbf{Head First Translation}: ``I don't have anything important to
  do right now, so let someone else use the CPU.''
\item
  \textbf{Godhood Tip}: Often used in lock-free programming spin-loops.
  If a lock-free CAS fails, you \passthrough{\lstinline!yield()!} to let
  the thread holding the lock finish its work faster.
\end{itemize}

\hypertarget{u.2-mutex-and-shared_mutex}{%
\subsection{\texorpdfstring{U.2 \texttt{\textless{}mutex\textgreater{}}
and
\texttt{\textless{}shared\_mutex\textgreater{}}}{U.2 \textless mutex\textgreater{} and \textless shared\_mutex\textgreater{}}}\label{u.2-mutex-and-shared_mutex}}

\hypertarget{stdtry_lock}{%
\subsubsection{\texorpdfstring{\texttt{std::try\_lock()}}{std::try\_lock()}}\label{stdtry_lock}}

\begin{itemize}
\tightlist
\item
  \textbf{Cppreference says}: Tries to lock the mutex. Returns
  immediately. On successful lock acquisition returns true, otherwise
  returns false.
\item
  \textbf{Head First Translation}: ``Is the bathroom door locked? If
  yes, I won't wait. I'll go do something else and come back later.''
\item
  \textbf{Godhood Tip}: This is a non-blocking operation. It is
  extremely useful in real-time systems (like games) where a thread
  cannot afford to block. If the mutex is locked, the thread abandons
  the task and moves on to the next frame.
\end{itemize}

\hypertarget{stdcall_once-and-stdonce_flag}{%
\subsubsection{\texorpdfstring{\texttt{std::call\_once} and
\texttt{std::once\_flag}}{std::call\_once and std::once\_flag}}\label{stdcall_once-and-stdonce_flag}}

\begin{itemize}
\tightlist
\item
  \textbf{Cppreference says}: Executes the Callable object exactly once,
  even if called concurrently, from several threads.
\item
  \textbf{Head First Translation}: ``The Ultimate Singleton Enforcer.''
\item
  \textbf{Godhood Tip}: This is the only thread-safe way to initialize
  global state or singletons before C++11's ``Magic Statics'' (where
  static local variables are thread-safe initialized).
\end{itemize}

\hypertarget{u.3-atomic}{%
\subsection{\texorpdfstring{U.3
\texttt{\textless{}atomic\textgreater{}}}{U.3 \textless atomic\textgreater{}}}\label{u.3-atomic}}

\hypertarget{stdatomicfetch_add-vs-stdatomicoperator}{%
\subsubsection{\texorpdfstring{\texttt{std::atomic::fetch\_add} vs
\texttt{std::atomic::operator++}}{std::atomic::fetch\_add vs std::atomic::operator++}}\label{stdatomicfetch_add-vs-stdatomicoperator}}

\begin{itemize}
\tightlist
\item
  \textbf{Cppreference says}: Atomically adds arg to the current value
  of the atomic object and returns the value held previously.
\item
  \textbf{Head First Translation}: ``Add 1 to the counter safely, but
  give me the number \emph{before} you added 1.''
\item
  \textbf{Godhood Tip}: \passthrough{\lstinline!fetch\_add!} returns the
  old value. If you need the new value, you have to add 1 to the result
  of \passthrough{\lstinline!fetch\_add!}, or just use
  \passthrough{\lstinline!operator++()!}. However,
  \passthrough{\lstinline!fetch\_add!} allows you to specify the
  \passthrough{\lstinline!memory\_order!}, whereas
  \passthrough{\lstinline!operator++!} always uses the heavy
  \passthrough{\lstinline!memory\_order\_seq\_cst!}. In high performance
  code, ALWAYS use
  \passthrough{\lstinline!fetch\_add(1, std::memory\_order\_relaxed)!}.
\end{itemize}

\hypertarget{stdatomiccompare_exchange_weak-vs-strong}{%
\subsubsection{\texorpdfstring{\texttt{std::atomic::compare\_exchange\_weak}
vs
\texttt{strong}}{std::atomic::compare\_exchange\_weak vs strong}}\label{stdatomiccompare_exchange_weak-vs-strong}}

\begin{itemize}
\tightlist
\item
  \textbf{Cppreference says}: Atomically compares the value
  representation of \passthrough{\lstinline!*this!} with that of
  \passthrough{\lstinline!expected!}. If they are bitwise-equal,
  replaces the former with \passthrough{\lstinline!desired!}.
\item
  \textbf{Head First Translation}: The CAS loop. We discussed this in
  Chapter 111.
\item
  \textbf{The Difference}: \passthrough{\lstinline!weak!} can fail
  ``spuriously'' (even if the values match, it might fail due to
  hardware reasons like a cache line eviction). You MUST put
  \passthrough{\lstinline!weak!} inside a
  \passthrough{\lstinline!while!} loop. \passthrough{\lstinline!strong!}
  will never fail spuriously, but it takes more CPU cycles.
\item
  \textbf{Godhood Tip}: If your algorithm requires a loop anyway (like
  traversing a linked list), use \passthrough{\lstinline!weak!}. If you
  don't have a loop, use \passthrough{\lstinline!strong!}.
\end{itemize}

\begin{center}\rule{0.5\linewidth}{0.5pt}\end{center}
